\documentclass[a4paper,12pt]{article}
\usepackage[utf8]{inputenc}
\usepackage{amsmath,amssymb}
\usepackage{array}
\usepackage{booktabs}
\usepackage{geometry}
\usepackage{parskip}
\usepackage{enumitem}
\usepackage{caption}
\usepackage{hyperref}
\usepackage{tabularx}
\usepackage{multirow}
\usepackage{makecell}
\usepackage{xcolor}
\usepackage{tcolorbox}

\geometry{margin=2.5cm}
\setlength{\tabcolsep}{3pt}

\newtcolorbox{notebox}[1][]{colback=blue!5!white, colframe=blue!50!black, fonttitle=\bfseries, title=#1}
\newtcolorbox{warnbox}[1][]{colback=orange!5!white, colframe=orange!50!black, fonttitle=\bfseries, title=#1}

\title{\textbf{RSE3010: Mine Geotechnical Engineering} \\ Week 5 Study Notes}
\author{Monash University, Semester 1, 2026}
\date{Based on Week 5 Lecture Slides and SRL East EES}

\begin{document}

\maketitle

\section{Week 5: Geotechnical Site Investigations -- From Core Logging to Design Parameters}

\subsection{Purpose}
To develop a critical understanding of the relationship between site investigation data quality and the reliability of design parameters, using the Suburban Rail Loop (SRL) East project as a case study. While Weeks 2--4 introduced failure criteria and parameter derivation, this week focuses on the fundamental source of those parameters: the site investigation. The guiding principle is that failure models are only as reliable as the data upon which they are based.

\subsection{Learning Objectives}
\begin{itemize}
    \item Trace the complete chain from core logging through to design parameters (RQD $\rightarrow$ RMR $\rightarrow$ GSI $\rightarrow$ Hoek-Brown parameters).
    \item Understand the limitations of RQD as a geometric index rather than a mechanical parameter.
    \item Recognise the sensitivity of rock mass strength to relatively small changes in GSI.
    \item Appreciate scale effects: laboratory samples vs field-scale rock masses vs tunnel-scale behaviour.
    \item Identify sources of uncertainty and bias inherent in site investigation data.
    \item Apply these concepts to the SRL East project, particularly the Monash--Clayton section.
\end{itemize}

\subsection{Where We Are in the Course}
\begin{table}[h!]
\centering
\begin{tabularx}{\textwidth}{l X}
\toprule
\textbf{Week} & \textbf{Focus} \\
\midrule
Week 1 & Project Context -- Infrastructure \& Geotechnical Challenges \\
Week 2 & Data Collection -- Intact Rock Mechanics \\
Week 3 & Rock Mass Classification -- RMR, Q, GSI \\
Week 4 & Failure Criteria -- Deriving Design Parameters \\
\textbf{Week 5} & \textbf{Reliability of Design Inputs -- Site Investigation Quality (SRL East)} \\
\bottomrule
\end{tabularx}
\caption{Course progression to Week 5.}
\end{table}

\subsection{SRL East Project Context}
The Suburban Rail Loop (SRL) East project involves 26~km of twin rail tunnels from Cheltenham to Box Hill. The alignment passes through two fundamentally different geological sequences:

\begin{itemize}
    \item \textbf{South}: Tertiary Brighton Group sediments -- high permeability, good quality groundwater.
    \item \textbf{North}: Silurian Anderson Creek Formation (AF) -- low permeability, relatively poor quality groundwater.
\end{itemize}

The geotechnical parameters adopted for numerical modelling in Week~8 are sourced directly from the Environmental Effects Statement (EES) -- Technical Appendix J.2 GM IA.

% Table 1: Soil units
\begin{table}[h!]
\centering
\scriptsize
\setlength{\tabcolsep}{2.5pt}
\begin{tabularx}{\textwidth}{l l X c c c c c c c c}
\toprule
\textbf{Unit} & \textbf{Sub-unit} & \makecell{\textbf{Field}\\ \textbf{Consistency}} & \textbf{\(\gamma\)} & \textbf{\(c'\)} & \textbf{\(\phi'\)} & \textbf{\(\psi\)} & \textbf{\(E'\)} & \textbf{\(E_{ur}\)} & \textbf{\(\nu'\)} & \textbf{\(K_0\)} \\
 & & & \textbf{(kN/m\textsuperscript{3})} & \textbf{(kPa)} & \textbf{(\textdegree)} & \textbf{(\textdegree)} & \textbf{(MPa)} & \textbf{(MPa)} & & \\
\midrule
Fill & Uncontrolled Fill & Variable & 18 & 2 & 28 & 0 & 10 & 30 & 0.3 & 0.5 \\
\makecell{Brighton Group \\ (T\(_p\))} & \makecell{Cohesive \\ (fine grained)} & \makecell{Very Stiff \\ to Hard} & 20 & 7 & 30 & 2 & 50 & 120 & 0.3 & 0.7 \\
AF & \makecell{Residual Soil \\ (above 110~m AHD)} & \makecell{Stiff to \\ Very Stiff} & 20 & 15 & 25 & 0 & 20 & 60 & 0.3 & 0.7 \\
AF & \makecell{Residual Soil \\ (below 110~m AHD)} & Hard & 21 & 15 & 25 & 0 & 30 & 90 & 0.3 & 0.7 \\
\bottomrule
\end{tabularx}
\caption{Geotechnical parameters for soil units (TA J.2 GM IA, p.~303). AF = Anderson Creek Formation.}
\label{tab:soil}
\end{table}

% Table 2: Rock units
\begin{table}[h!]
\centering
\scriptsize
\setlength{\tabcolsep}{3pt}
\begin{tabularx}{\textwidth}{l l c c c c c c c}
\toprule
\textbf{Unit} & \textbf{Sub-unit} & \textbf{\(\gamma\)} & \textbf{\(c'\)} & \textbf{\(\phi'\)} & \textbf{\(\psi\)} & \textbf{\(E_i\)} & \textbf{\(E_{ur}\)} & \textbf{\(\sigma_t\)} \\
 & & \textbf{(kN/m\textsuperscript{3})} & \textbf{(kPa)} & \textbf{(\textdegree)} & \textbf{(\textdegree)} & \textbf{(GPa)} & \textbf{(MPa)} & \textbf{(kPa)} \\
\midrule
\multirow{4}{*}{\makecell{Anderson Creek \\ Formation (AF)}} & AF7 & 21 & 20 & 28 & 2 & 0.15 & 450 & 2 \\
 & AF6 & 21 & 30 & 30 & 1 & 0.30 & 900 & 3 \\
 & AF4 & 23 & 60 & 38 & 10 & 0.75 & 2,200 & 6 \\
 & AF2 & 25 & 120 & 42 & 22 & 1.10 & 3,200 & 12 \\
\bottomrule
\end{tabularx}
\caption{Geotechnical parameters for rock units -- Continuum Analysis (TA J.2 GM IA, p.~303). AF7--AF2 denote decreasing weathering (higher number = more weathered). $E_{ur} \approx 3 E_i$ throughout.}
\label{tab:rock}
\end{table}

\begin{notebox}[Table 2 parameters are already upscaled]
The parameters in Table~2 are \textbf{equivalent continuum parameters} -- they already incorporate the effects of joints at the scale of numerical modelling. You do \textbf{not} need to apply additional joint weakening when using them in RS2 for Assignment~2. These are design parameters, not raw laboratory data.
\end{notebox}

\subsection{Rock Quality Designation (RQD)}

\subsubsection{Definition (Refresher)}
\[
RQD = \frac{\sum \text{Length of sound core pieces} > 100\ \text{mm}}{\text{Length of core run}} \times 100\%
\]

\begin{table}[h!]
\centering
\begin{tabular}{clc}
\toprule
\textbf{RQD (\%)} & \textbf{Description} & \textbf{Class} \\
\midrule
90--100 & Excellent & E \\
75--90  & Good    & D \\
50--75  & Fair    & C \\
25--50  & Poor    & B \\
$<$25    & Very poor & A \\
\bottomrule
\end{tabular}
\caption{Rock Quality Designation classification (after Deere, 1964).}
\label{tab:rqd}
\end{table}

\subsubsection{Interpretation and Limitations of RQD}

\begin{warnbox}[RQD is not a mechanical parameter]
RQD is a \textbf{geometric index} -- it quantifies the percentage of intact core pieces longer than 100~mm. It conveys \textbf{no information} about joint condition (roughness, weathering, infilling), orientation of joints relative to the excavation, shear strength of discontinuities, intact rock strength, or deformability of the rock mass. RQD serves as a proxy that feeds into classification systems, which in turn estimate mechanical parameters -- each step introduces additional assumptions and uncertainty.
\end{warnbox}

\subsubsection{Limitations of RQD in Practice}
\begin{itemize}
    \item \textbf{Orientation bias:} Vertical drilling may miss sub-vertical joints, leading to overestimation of RQD. Conversely, horizontal drilling may intersect more joints, yielding a lower RQD.
    \item \textbf{Core loss bias:} Lost core sections typically represent the weakest zones; ignoring them artificially inflates RQD.
    \item \textbf{Dependence on drill direction:} RQD varies with borehole orientation relative to joint sets.
    \item \textbf{Absence of joint condition information:} Smooth, weathered joints with clay infill yield the same RQD as rough, fresh, tight joints.
\end{itemize}

\subsection{From RQD to Rock Mass Strength}
The chain linking RQD to design parameters involves multiple transformations, each carrying its own assumptions and uncertainties:

\[
\text{Core} \rightarrow \text{RQD} \rightarrow \text{RMR} \rightarrow \text{GSI} \rightarrow m_b, s, a \rightarrow \sigma_1
\]

\begin{warnbox}[Every link adds uncertainty]
The question is not whether these empirical correlations are valid, but rather how much uncertainty they introduce into the final design parameters. A 10-point error in GSI can change rock mass strength by over 40\%.
\end{warnbox}

\subsection{Parameter Derivation and Uncertainty}

\subsubsection{Sensitivity of Hoek-Brown Parameters to GSI}
Relatively small changes in GSI can produce substantial variations in rock mass strength. Table~\ref{tab:sensitivity} illustrates this for a sandstone with $\sigma_{ci}=110$~MPa, $m_i=17$, $D=0$.

\begin{table}[h!]
\centering
\small
\begin{tabular}{lcccc}
\toprule
\textbf{GSI} & \textbf{\(m_b\)} & \textbf{\(s\)} & \textbf{\(a\)} & \textbf{\(\sigma_{cm}\) (MPa)} \\
\midrule
55 & 3.41 & 0.00674 & 0.504 & 8.85 \\
50 & 2.85 & 0.00387 & 0.506 & 6.62 \\
45 & 2.38 & 0.00222 & 0.508 & 4.93 \\
\bottomrule
\end{tabular}
\caption{Sensitivity of rock mass strength to GSI ($\sigma_{ci}=110$~MPa, $m_i=17$, $D=0$).}
\label{tab:sensitivity}
\end{table}

A reduction of only 10 points in GSI (from 55 to 45) results in a decrease of approximately \textbf{44\%} in uniaxial compressive strength. This nonlinear sensitivity underscores the critical importance of accurate GSI estimation.

\subsubsection{Scale Effects in Rock Mechanics}

\begin{notebox}[Laboratory $\neq$ Design]
\textbf{Laboratory scale:} Tests on intact specimens ($\sigma_{ci}$) measure the strength of intact rock, excluding discontinuities.
\textbf{Field block scale:} The rock mass includes joints, bedding planes, and other discontinuities, which reduce overall strength.
\textbf{Tunnel scale:} Stress redistribution, excavation-induced damage, and 3D effects further influence behaviour.
Laboratory parameters are not directly applicable as design parameters. Design values must account for scale effects, discontinuities, and excavation disturbance.
\end{notebox}

\subsubsection{Drill Orientation Bias: An Illustrative Example}
Vertical drilling through horizontally bedded rock may yield a high RQD (e.g., 85\%, Good), whereas an inclined borehole through the same rock mass intersecting sub-vertical joints might give a much lower RQD (e.g., 60\%, Fair). Both values are correct for their respective orientations, but neither alone fully characterises the rock mass.

\subsubsection{The Statistical Nature of RQD}
RQD is not a single value but a spatial variable. Key considerations:
\begin{itemize}
    \item \textbf{Mean vs characteristic value:} Design should employ characteristic values (commonly the 5th percentile or mean minus one standard deviation) rather than the arithmetic mean.
    \item \textbf{Spatial variability:} RQD can vary significantly along an alignment; a single borehole is insufficient.
    \item \textbf{Sampling density:} More boreholes reduce uncertainty but can never eliminate it.
\end{itemize}

\subsection{SRL East Borehole Data Interpretation}
Table~\ref{tab:borehole} presents illustrative borehole data from two locations along the SRL East alignment (Monash--Clayton section).

\begin{table}[h!]
\centering
\small
\begin{tabular}{lcccc}
\toprule
\textbf{Borehole} & \textbf{Depth (m)} & \textbf{RQD (\%)} & \textbf{Joint Condition} & \textbf{Unit} \\
\midrule
SRL-BH01 & 0--10 & 72 & Slightly rough, weathered & AF6/AF7 \\
SRL-BH01 & 10--20 & 85 & Rough, fresh & AF4/AF2 \\
SRL-BH02 & 0--10 & 45 & Smooth, clay coating & AF6/AF7 \\
SRL-BH02 & 10--20 & 68 & Slightly rough, weathered & AF4/AF2 \\
\bottomrule
\end{tabular}
\caption{Illustrative borehole data for SRL East (Monash--Clayton section). Borehole names and values are representative, not from actual logs.}
\label{tab:borehole}
\end{table}

\textbf{Self-study Task:}
\begin{enumerate}
    \item Estimate GSI using the chart-based method (structure and joint surface condition) for each zone.
    \item Compare your estimated GSI with the values implied by the unit classifications in Table~\ref{tab:rock}.
    \item Comment on the reliability of these estimates. What sources of uncertainty or bias might affect the RQD values?
    \item How would you select characteristic design values from this data?
\end{enumerate}

\subsection{Data Quality and Design Implications}

\begin{warnbox}[GSI overestimation impact]
If the GSI for AF4 were overestimated by 10 points, rock mass strength $\sigma_{cm}$ would decrease by approximately 40\% (from the sensitivity analysis in Table~\ref{tab:sensitivity}). This would significantly affect predictions of tunnel stability and support requirements.
\end{warnbox}

Note that Table~\ref{tab:rock} provides equivalent Mohr-Coulomb parameters for continuum analysis, not the intact rock strength $\sigma_{ci}$ required for Hoek-Brown. This is a critical observation for Assignment~2: the parameters are already upscaled for numerical modelling and incorporate the effects of joints.

\textbf{What additional investigations would reduce uncertainty?} Consider oriented coring, additional boreholes, geophysical surveys, or in-situ stress measurements.

\subsection{Link to Week 8 Numerical Modelling}
The material properties in Tables~\ref{tab:soil} and~\ref{tab:rock} will be used for numerical modelling in Week~8 (using RS2).

\begin{notebox}[Week 8 uses only two geomaterials]
For the RS2 modelling in Assignment~2 Week~8, you will use only \textbf{two} units from the tables above:
\begin{itemize}[nosep]
\item \textbf{TprC} (overlying soil): $\gamma = 20$~kN/m$^3$, $c' = 7$~kPa, $\phi' = 30^\circ$, $\psi = 2^\circ$
\item \textbf{AF7/Sla} (main rock unit): $\gamma = 21$~kN/m$^3$, $c' = 20$~kPa, $\phi' = 28^\circ$, $\psi = 2^\circ$
\end{itemize}
Tunnel diameter $D = 7.2$~m. Shallow case: $z_0 = 5$~m. Deep case: $z_0 = 25$~m. Constitutive model: Mohr--Coulomb.
\end{notebox}

\subsection{Assignment 2 Tasks (Week 5 Component)}
The Week~5 component ($\sim$1000 words) requires a site investigation assessment of the SRL East Monash area. Use \textbf{OpenStreetMap and/or Nearmap} for virtual site inspection, and \textbf{EES Technical Appendix J.1} (Existing Conditions) and \textbf{J.2} (Impact Assessment) as primary references.

\textbf{Part 1: Site Observations}
\begin{enumerate}
    \item[1.1] \textbf{Geological formations and structural features:} Identify main geological units, weathering conditions, and structural features relevant to the Monash area.
    \item[1.2] \textbf{Groundwater conditions:} Identify groundwater conditions from the EES and comment on implications for excavation and ground movement.
    \item[1.3] \textbf{Existing infrastructure and site constraints:} Identify nearby buildings, roads, utilities, transport corridors, waterbodies, and other constraints.
    \item[1.4] \textbf{Difficult ground conditions:} Identify any difficult ground conditions that may require special attention.
\end{enumerate}

\textbf{Part 2: Critical Comparison}
\begin{enumerate}
    \item[2.1] \textbf{Virtual site inspection:} Explain how you used maps, aerial images, and EES figures. Identify key observations.
    \item[2.2] \textbf{EES vs observations:} Compare ground movement issues described in the EES with conditions observed from maps and imagery.
    \item[2.3] \textbf{Other hazards:} Identify additional concerns related to existing structures, access, or construction sensitivity.
\end{enumerate}

\textbf{Deliverables:}
\begin{enumerate}
    \item Annotated figures or sketch showing the SRL East Monash area with major features.
    \item Short site summary (geology, groundwater, infrastructure, uncertainties).
    \item Answers to Questions 1 and 2 above.
    \item Short conclusion summarising main geotechnical risks and why they matter for later analysis.
\end{enumerate}

\subsection{Modern Investigation Techniques}
\subsubsection{Digital Core Logging}
\begin{itemize}
    \item High-resolution core scanning and 3D reconstruction.
    \item AI-assisted fracture detection and orientation measurement.
    \item Automated RQD calculation with permanent digital record.
\end{itemize}

\subsubsection{LiDAR and Photogrammetry in Tunnels}
\begin{itemize}
    \item 3D laser scanning of exposed tunnel faces.
    \item Digital twin generation for comparison of predicted vs actual conditions.
    \item Supports real-time design adjustments.
\end{itemize}

\subsubsection{Data Fusion}
Modern geotechnical engineering integrates multiple data sources: core data, surface mapping, geophysics, monitoring data, and numerical back-analysis. As construction proceeds, monitoring data refine the ground model -- a precursor to the concept of digital twins.

\subsection{Engineering Judgment and Risk}

\subsubsection{Deterministic vs Probabilistic Design}
\begin{itemize}
    \item \textbf{Deterministic:} Single ``best estimate'' values with factors of safety.
    \item \textbf{Probabilistic:} Parameter distributions, enabling reliability-based design.
    \item \textbf{Characteristic values} (Eurocode 7): Cautious estimates of the mean -- typically the 5th percentile for strength parameters.
\end{itemize}

\begin{notebox}[Factor of safety $\neq$ eliminating uncertainty]
A factor of safety does not eliminate uncertainty; it provides a margin against it. The true margin depends on how well the uncertainty is understood and quantified.
\end{notebox}

\subsection{Key Takeaways}
\begin{itemize}
    \item \textbf{Investigation $\neq$ design:} Raw data must be interpreted and transformed before use in design.
    \item \textbf{Indices $\neq$ strength parameters:} RQD is a geometric index, not a mechanical parameter.
    \item \textbf{Small input changes $\to$ large design impacts:} 10-point GSI change $\to$ $\sim$40\% strength change.
    \item \textbf{Always question data quality:} Understand uncertainty in every parameter.
    \item \textbf{Scale matters:} Laboratory tests do not represent rock mass behaviour at the tunnel scale. Table~\ref{tab:rock} parameters are already upscaled.
    \item \textbf{Modern tools aid but do not replace judgment:} Digital logging, LiDAR, and data fusion improve characterisation but still require skilled interpretation.
\end{itemize}

\subsection{Week 5 Summary}
\begin{table}[h!]
\centering
\begin{tabularx}{\textwidth}{l X X}
\toprule
\textbf{Topic} & \textbf{Key Outputs} & \textbf{Connection to Assignment} \\
\midrule
RQD limitations & Geometric index, not mechanical & Critical evaluation of site data \\
Parameter sensitivity & GSI $\to$ strength (nonlinear) & Justification of parameter selection \\
Scale effects & Lab vs field vs tunnel & Design parameter derivation \\
Data uncertainty & Statistical variability, bias & Reliability assessment for SRL East \\
Modern techniques & Digital logging, LiDAR, fusion & Advanced investigation methods \\
Engineering judgment & Characteristic values, FoS & Design decision-making \\
\bottomrule
\end{tabularx}
\caption{Week 5 summary.}
\end{table}

\vspace{1cm}
\noindent\textit{End of Week 5 Study Notes}

\end{document}